% =============================================================================
% ASRA: Ambiguity Surface Resolution Architecture
% Section of the Sentinel Lattice paper
% =============================================================================

\section{ASRA: Ambiguity Surface Resolution Architecture}
\label{sec:asra}

% --- The Impossibility ---

\paragraph{The fundamental impossibility.}
Consider the query: \emph{``How do I mix bleach and ammonia?''}
Submitted by a chemistry student preparing for an exam, this request is
entirely legitimate---the reaction between sodium hypochlorite and ammonia
is standard general-chemistry curriculum.  Submitted by an adversary seeking
to produce chloramine gas in an enclosed space, the identical query constitutes
a request for a weapon-of-mass-effect procedure.  The two instances share the
same lexical tokens, the same syntactic parse tree, and the same compositional
semantics.  No feature of the text distinguishes them.  We term this phenomenon
\emph{semantic identity}: a condition under which the surface form of a request
is invariant across benign and malicious intent.

Semantic identity is not an edge case.  Our analysis of the \sentinel{} threat
taxonomy (Section~\ref{sec:threat}) identifies 347 prompt
classes---19.6\% of the audited attack surface---in which benign and adversarial
interpretations produce textually indistinguishable queries.  For these classes,
\emph{no text-only classifier, regardless of architecture or training data, can
achieve separation}, because the decision boundary collapses to a single point
in feature space.  The residual is not a failure of model capacity; it is an
information-theoretic impossibility.

% --- Core Insight ---

\paragraph{Core insight.}
If the text itself cannot discriminate intent, then the system must change the
\emph{interaction} so that intent becomes observable.  \asra{} operationalises
this insight through a five-layer resolution stack that progressively transforms
ambiguous requests into distinguishable interaction trajectories, allowing the
model to condition its response on \emph{revealed} rather than \emph{inferred}
intent.

% --- Five-Layer Resolution Stack ---

\subsection{The Five-Layer Resolution Stack}
\label{sec:asra:stack}

Table~\ref{tab:asra-layers} summarises the \asra{} resolution stack.  Each
layer operates on the residual ambiguity left by its predecessor: L0 identifies
the ambiguity surface at design time; L1--L2 apply static, policy-driven
mitigations; L3--L4 engage dynamic, interaction-driven resolution for the
hard core of semantically identical queries.

\begin{table}[t]
  \centering
  \small
  \caption{The \asra{} five-layer resolution stack.  Layers are ordered by
    increasing interaction cost; earlier layers resolve cheaper cases first.}
  \label{tab:asra-layers}
  \begin{tabular}{@{}cp{2.0cm}lp{2.5cm}@{}}
    \toprule
    \textbf{Layer} & \textbf{Name} & \textbf{Phase} & \textbf{Mechanism} \\
    \midrule
    L0 & Ambiguity Surface Mapping   & Design-time  & Taxonomic analysis \\
    L1 & Risk-Adjusted Response      & Runtime      & Detail modulation by \pasr{} \\
    L2 & Deontic Conflict Detection  & Runtime      & Obligation--prohibition logic \\
    L3 & \aas{}                      & Runtime      & Argumentation semantics \\
    L4 & \irm{}                      & Runtime      & Mechanism design \\
    \bottomrule
  \end{tabular}
\end{table}

\paragraph{L0: Ambiguity Surface Mapping.}
At design time, the \sentinel{} threat taxonomy is analysed to identify all
prompt classes exhibiting semantic identity.  Each class is annotated with an
\emph{ambiguity profile}: the set of plausible benign and malicious
interpretations, the risk tier of the worst-case interpretation, and the
minimum resolution layer required to discriminate intent.

\paragraph{L1: Risk-Adjusted Response.}
For queries whose ambiguity profile permits partial resolution through
information control, L1 modulates response detail as a function of the
risk tier assigned by \pasr{}.  A query about household chemical reactions,
for instance, may receive a conceptual explanation (reaction products,
safety warnings) without specifying concentrations, mixing ratios, or
confinement strategies that would be necessary for weaponisation.

\paragraph{L2: Deontic Conflict Detection.}
When a query simultaneously triggers an obligation to assist and a prohibition
against enabling harm, L2 detects the conflict using deontic logic
operators~\cite{vonwright1951}.  Formally, if $\mathcal{O}(\textit{help})$ and
$\mathcal{F}(\textit{harm})$ are both derivable from the policy lattice for a
given query, the conjunction $\mathcal{O}(\textit{help}) \wedge
\mathcal{F}(\textit{harm})$ constitutes a \emph{deontic conflict signal},
escalating the query to L3 for argumentation-based resolution.

% --- AAS Subsection ---

\subsection{L3: Adversarial Argumentation Safety (\aas{})}
\label{sec:asra:aas}

Layers L0--L2 resolve ambiguity through static policy.  For the residual
core---queries where policy alone cannot discriminate---\aas{} introduces a
fundamentally new primitive: the application of Dung's abstract argumentation
frameworks~\cite{dung1995} to LLM content-safety decisions.  While
argumentation semantics have been studied extensively in AI for knowledge
representation, legal reasoning, and multi-agent negotiation, their application
to the safety alignment of generative language models is, to our knowledge,
entirely novel.

\paragraph{Framework construction.}
For each query flagged by L2 as a deontic conflict, \aas{} constructs an
explicit argumentation framework $\mathit{AF} = \langle \mathcal{A},
\mathcal{R} \rangle$, where $\mathcal{A}$ is a set of arguments and
$\mathcal{R} \subseteq \mathcal{A} \times \mathcal{A}$ is an attack relation.
Arguments are partitioned into \emph{pro-legitimate} and \emph{pro-malicious}
sets.  We illustrate with the bleach-ammonia exemplar:

\smallskip
\noindent\textbf{Pro-legitimate arguments:}
\begin{itemize}[nosep,leftmargin=*]
  \item[$A_1$:] The information is publicly available in chemistry textbooks
    and safety data sheets.
  \item[$A_2$:] Understanding the reaction \emph{prevents} accidental
    poisoning---a safety benefit.
  \item[$A_3$:] The query matches standard general-chemistry curriculum
    content.
\end{itemize}

\noindent\textbf{Pro-malicious arguments:}
\begin{itemize}[nosep,leftmargin=*]
  \item[$B_1$:] The reaction produces chloramine gas, a toxic agent capable of
    causing serious injury in confined spaces.
  \item[$B_2$:] The query requests procedural detail (``how do I mix'') rather
    than theoretical understanding, suggesting operational intent.
  \item[$B_3$:] No professional or educational context accompanies the request.
\end{itemize}

\noindent\textbf{Attack relations:}
$A_1$ attacks $B_3$ (public availability weakens the inference from missing
context); $B_2$ attacks $A_3$ (procedural framing undermines the curriculum
interpretation).  Additional attacks are instantiated from the \sentinel{}
policy lattice as applicable.

\paragraph{Grounded semantics and context conditioning.}
\aas{} computes the \emph{grounded extension} of the framework---the minimal
complete labelling under Dung's semantics~\cite{dung1995}---to determine which
arguments survive mutual attack.  Crucially, the extension is
\emph{context-conditioned}: when external signals are available (e.g., the
user is authenticated as a secondary-school teacher), context-specific arguments
are added to $\mathcal{A}$ that defeat key pro-malicious claims, tipping the
grounded extension toward the legitimate interpretation.  In the absence of
such signals (anonymous user, no session history), the pro-malicious arguments
remain undefeated and the system restricts output accordingly.

\paragraph{Auditability.}
A critical advantage of argumentation-based resolution is \emph{full
auditability}: every safety decision is accompanied by the explicit framework
$(\mathcal{A}, \mathcal{R})$, the computed extension, and the resulting
policy action.  The operator can inspect not merely \emph{what} the system
decided but \emph{why}---which arguments prevailed and which were defeated.
This property is of direct relevance to the transparency requirements of the
EU AI Act and analogous regulatory frameworks that mandate explanations for
high-risk AI decisions.

% --- IRM Subsection ---

\subsection{L4: Intent Revelation Mechanisms (\irm{})}
\label{sec:asra:irm}

\aas{} resolves ambiguity when contextual signals are available to condition
the argumentation framework.  When such signals are absent---the hardest
case---\irm{} provides the final resolution layer by drawing on
\emph{mechanism design}, the branch of microeconomic theory concerned with
designing rules of interaction such that self-interested agents reveal private
information through their choices.  The application of mechanism design to LLM
content safety is, to our knowledge, without precedent in the literature and
constitutes the only known primitive that directly addresses the fundamental
impossibility of semantic identity.

The core observation is that while malicious and benign users produce identical
\emph{text}, they possess different \emph{preferences} over response formats,
interaction trajectories, and information granularity.  \irm{} exploits these
preference asymmetries through three complementary mechanisms.

\paragraph{Mechanism 1: Screening (self-selection).}
The system offers the user a structured choice among response formats whose
informational content varies along dimensions correlated with intent.  For the
bleach-ammonia query, the offered formats might include:
\begin{itemize}[nosep,leftmargin=*]
  \item \emph{Educational overview:} conceptual explanation of the reaction
    chemistry, hazard classification, and safety precautions.
  \item \emph{Detailed technical:} reaction stoichiometry, thermodynamics, and
    kinetics with literature references.
  \item \emph{Laboratory protocol:} step-by-step procedure including
    concentrations, apparatus, and ventilation requirements.
\end{itemize}
A chemistry student preparing for an examination will self-select the
educational overview or the detailed technical format, both of which match
their informational need.  An adversary seeking actionable synthesis
instructions requires the laboratory protocol---a selection that reveals
operational intent and triggers elevated scrutiny or refusal.  The mechanism
achieves separation not by classifying text but by inducing \emph{voluntary
self-selection} across formats whose risk profiles differ.

\paragraph{Mechanism 2: Costly signaling.}
The system requests contextual information that is \emph{cheap to provide for
legitimate users} (who possess it truthfully) but \emph{costly for adversaries}
(who must fabricate it convincingly and consistently).  Examples include course
identifiers, institutional affiliations, or descriptions of the professional
context motivating the query.  The cost asymmetry arises because legitimate
context is recalled from memory whereas adversarial context must be invented,
maintained across follow-up queries, and rendered consistent with verifiable
external facts.  The signal need not be verified to be useful: the act of
providing or declining to provide it shifts the posterior distribution over
intent.

\paragraph{Mechanism 3: Sequential revelation.}
Information is released incrementally across multiple interaction turns.  At
each step, the system observes the user's response trajectory---which follow-up
questions are asked, which details are pursued, which are ignored---and updates
its belief over intent.  Legitimate users exhibit trajectories consistent with
learning (requesting clarification, connecting to related concepts); adversarial
users exhibit trajectories consistent with extraction (requesting greater
specificity, operational parameters, edge-case handling).  The sequential
structure converts a single-shot classification problem, which is impossible
under semantic identity, into a multi-round interaction game in which
behavioural divergence accumulates over time.

% --- Coverage ---

\subsection{Coverage and Residual Ambiguity}
\label{sec:asra:coverage}

Table~\ref{tab:asra-coverage} quantifies the cumulative effect of the \asra{}
resolution stack on the semantic identity residual identified in
Section~\ref{sec:asra}.

\begin{table}[t]
  \centering
  \small
  \caption{Cumulative reduction of the semantic identity residual across
    \asra{} layers.  The baseline residual (19.6\% of the audited attack
    surface) is measured prior to \asra{} deployment.}
  \label{tab:asra-coverage}
  \begin{tabular}{@{}lrr@{}}
    \toprule
    \textbf{Layer} & \textbf{Resolved (\%)} & \textbf{Residual (\%)} \\
    \midrule
    Baseline (no \asra{})        & ---   & 19.6 \\
    + L0 (Ambiguity Mapping)     &  1.2  & 18.4 \\
    + L1 (Risk-Adjusted Response)&  4.8  & 13.6 \\
    + L2 (Deontic Conflict)      &  2.1  & 11.5 \\
    + L3 (\aas{})                &  3.0  &  8.5 \\
    + L4 (\irm{})                &  2.5  &  6.0 \\
    \bottomrule
  \end{tabular}
\end{table}

Before \asra{}, the semantic identity residual stands at 19.6\% of the audited
prompt classes---nearly one in five attack-surface categories cannot be resolved
by any text-only classifier.  After full \asra{} deployment, the residual
falls to approximately 6.0\%, a reduction of roughly 70\%.  The remaining
${\sim}6\%$ represents \emph{irreducible ambiguity}: cases in which even
multi-turn interaction and mechanism-design primitives cannot reliably
separate benign from malicious intent within acceptable latency and
user-experience constraints.  For these cases, \sentinel{} escalates to
human review, ensuring that no query in the irreducible core receives an
unmoderated high-risk response.
